\section{Prova anterior: primeira prova teórica}

Esta seção reproduz os enunciados visíveis de uma prova anterior, para uso como
exercício de revisão. Não são apresentadas respostas.

\subsection{Questão 1 (1,0 ponto)}

Um(a) geólogo(a) foi a campo numa jornada de 15 dias para mapear uma área
sedimentar de \SI{15}{\kilo\metre\squared}. Ao longo dos dias, obteve um total de
150 pontos. Desses 150 pontos, o(a) geólogo(a) sumarizou em 20 as fácies
sedimentares encontradas. Dessas 20, teve dúvidas acerca dos mecanismos,
processos e tipos de fluxos capazes de gerar as seguintes fácies:

\begin{description}[style=nextline,leftmargin=1.5cm,labelsep=0.5em]
\item[\textbf{Fácies A}] Arenito cinza, com granulometria de areia fina a
média, bem selecionado, aparentemente com alta maturidade mineralógica e
textural e com estrutura maciça.

\begin{itemize}
    \item A cor cinza pode estar relacionada à presença de minerais escuros,
    matéria orgânica ou alterações diagenéticas.

    \item A boa seleção indica que grãos de tamanhos semelhantes foram
    transportados seletivamente ou retrabalhados.

    \item A alta maturidade mineralógica indica o predomínio de minerais
    estáveis, principalmente quartzo. A maturidade textural está relacionada
    à pouca matriz, à boa seleção e ao maior arredondamento dos grãos.

    \item A estrutura maciça indica ausência de estruturas sedimentares
    internas visíveis. Isso pode resultar de deposição rápida ou da
    destruição das estruturas originais por bioturbação.

    \item O transporte provavelmente ocorreu por tração e saltação. A
    deposição ocorreu quando o fluxo perdeu competência. Como não existem
    estruturas internas visíveis, não é possível determinar com segurança
    se o fluxo era unidirecional ou oscilatório.
\end{itemize}

A fácies é compatível com um quartzoarenito bem selecionado e com pouca matriz.

  \item[\textbf{Fácies B}] Arenito intercalado com lamito, na proporção 1:1,
  com ondulações;

A proporção aproximadamente igual de areia e lama caracteriza um acamamento
ondulado do tipo \textit{wavy bedding}. Durante os períodos de maior energia,
a areia foi transportada por tração e saltação, formando ripples.

Com a redução da velocidade e da turbulência do fluxo, as partículas finas
mantidas em suspensão depositaram-se por decantação sobre os ripples,
acompanhando seu perfil. A coesão da lama posteriormente depositada favoreceu
sua preservação.

Com o retorno de um fluxo mais energético, a areia voltou a ser transportada
e depositada. A repetição entre períodos de corrente e períodos de baixa
energia formou a intercalação entre arenito e lamito.

Essa alternância é compatível com variações periódicas da energia do fluxo,
como as que podem ocorrer sob influência de marés, embora a fácies isolada
não seja suficiente para determinar o ambiente deposicional.

Esse é uma fácies compatível com uma área onde ocorria na corrente o transporte por tração e no estofo da maré havia decantação de finos.

\item[\textbf{Fácies C}] Folhelho cinza-escuro, com cristais de pirita e
presença do icnofóssil \textit{Zoophycos}.

\begin{itemize}
    \item Os folhelhos são rochas finas, formadas principalmente por partículas
    do tamanho silte e argila, e apresentam fissilidade.

    \item A cor cinza-escura pode indicar a presença e preservação de matéria
    orgânica em condições redutoras.

    \item A pirita, de composição $\mathrm{FeS_2}$, forma-se pela combinação
    entre ferro e enxofre, geralmente durante a diagênese e em condições
    pobres em oxigênio.

    \item O icnofóssil \textit{Zoophycos} registra bioturbação, alimentação e
    retrabalhamento do sedimento por organismos.

    \item As partículas finas foram transportadas em suspensão e depositadas
    por decantação em condições de baixa energia.
\end{itemize}

A presença de \textit{Zoophycos} indica que, apesar das condições redutoras,
existiu oxigênio suficiente em determinados períodos ou níveis do sedimento
para permitir atividade biológica.

\item[\textbf{Fácies D}] Arenito amarelado, de granulação média, com presença
de lineação e estrutura plano-paralela.

\begin{itemize}
    \item A granulação média é compatível com o transporte da areia como carga
    de fundo, principalmente por tração e saltação.

    \item A lineação é interpretada como uma lineação de corrente, produzida
    pelo alinhamento de grãos ou por marcas paralelas ao fluxo. Ela pode indicar
    a direção da paleocorrente.

    \item A estrutura plano-paralela indica transporte e deposição sobre um
    leito plano, sob ação de uma corrente unidirecional relativamente energética.

    \item A coloração amarelada pode estar relacionada à presença de óxidos ou
    hidróxidos de ferro formados em condições oxidantes.
\end{itemize}

A fácies é compatível com o transporte de areia por uma corrente unidirecional,
com tração e saltação sobre um leito plano.


\item[\textbf{Fácies E}] Arenito de coloração rosácea, com estratificação
cruzada de baixo ângulo e com o icnofóssil \textit{Skolithos}.

\begin{itemize}
    \item A coloração rosácea pode estar relacionada à presença de óxidos de
    ferro, principalmente hematita, formados em condições oxidantes.

    \item A estratificação cruzada de baixo ângulo registra a migração de formas
    de leito de baixo relevo. A areia foi transportada principalmente por tração
    e saltação.

    \item O icnofóssil \textit{Skolithos} corresponde a tubos de habitação
    aproximadamente verticais, produzidos por organismos que viviam em
    substrato arenoso e oxigenado.

    \item Esses organismos provavelmente se alimentavam de partículas
    suspensas na água; portanto, sua presença não exige um sedimento rico em
    matéria orgânica.
\end{itemize}

A fácies registra transporte de areia, migração de formas de leito e posterior
ou simultânea ocupação do substrato por organismos.

\end{description}

Com base nessas informações, esclareça os mecanismos, processos e tipos de
fluxos que geraram as fácies sedimentares citadas.

\begin{quote}
  \textit{Observação da prova: não é necessário realizar associação de fácies.}
\end{quote}

\subsection{Questão 2 (1,0 ponto)}

\begin{quote}
``Uma das novidades contidas na estratigrafia moderna é o uso das fácies
sedimentares, a partir de estudos de sucessões verticais, na arte de interpretar
ambientes deposicionais antigos. Esta é uma das tarefas mais árduas da Geologia
Sedimentar, por envolver normalmente preconceitos arraigados e presença de
incerteza. Por outro lado, é importante em termos econômicos, já que encerra os
elementos básicos para a exploração de riquezas minerais como os energéticos
(petróleo, carvão, urânio), além de outros depósitos valiosos (evaporitos e
\textit{placers} --- ouro, diamante e cassiterita, como exemplos).''

\hfill (Della Fávera, 2001, p.~121).
\end{quote}

Uma fácies corresponde a um corpo de sedimentos ou de rochas que apresenta
características próprias, como litologia, granulometria, estruturas
sedimentares, fósseis, geometria e tipos de contato. Essas características
permitem diferenciar esse corpo dos depósitos vizinhos e interpretar os
processos responsáveis por sua formação.

Entre os principais tipos estão as litofácies, definidas principalmente pelas
características físicas e litológicas; as biofácies, reconhecidas pelo conteúdo
fossilífero; e as icnofácies, caracterizadas pelos vestígios da atividade dos
organismos, como pistas, escavações e outros icnofósseis.

A descrição puramente litológica identifica a composição e as características
da rocha, como arenito, folhelho ou conglomerado. A análise de fácies utiliza
essas informações juntamente com as estruturas sedimentares, os fósseis, a
geometria e os contatos para interpretar os processos de transporte e
deposição.

Fácies que aparecem relacionadas podem ser reunidas em associações de fácies.
Sua organização lateral e vertical forma sucessões que ajudam a reconstruir as
mudanças ocorridas no ambiente deposicional.

Uma forma de compreender essa organização é considerar os rios meandrantes,
cujos canais mudam de posição ao longo do tempo geológico. Uma área
inicialmente ocupada pela calha do rio pode tornar-se um canal abandonado e,
posteriormente, ser recoberta por depósitos da planície de inundação. Assim,
os sedimentos arenosos do canal e das barras em pontal podem ficar recobertos
por sedimentos finos da planície.

Portanto, o estudo das fácies e de seu empilhamento permite reconhecer os
processos sedimentares e reconstruir os ambientes deposicionais antigos.